% exercises.tex
\documentclass[11pt, a4paper]{article}

\usepackage[T1]{fontenc}
\usepackage[utf8]{inputenc}
\usepackage{tgpagella}
\usepackage[spanish, es-noshorthands]{babel}
\usepackage{amsmath, amssymb, bm}
\usepackage[margin=2.5cm]{geometry}
\usepackage{fancyhdr}
\usepackage[most]{tcolorbox}

\pagestyle{fancy}
\fancyhf{}
\setlength{\headheight}{16pt}
\lhead{\textbf{UCB Sede Tarija} | Ing. Mecatrónica}
\rhead{Ejercicios Asignados: Subconjunto Fundacional[cite: 4]}
\lfoot{Prof: M.Sc. Ing. Bernardo Quiroga Turdera[cite: 1, 3]}
\renewcommand{\headrulewidth}{0.4pt}
\definecolor{ucbblue}{RGB}{0, 51, 102}

\begin{document}

\begin{tcolorbox}[colback=ucbblue!5!white, colframe=ucbblue, title=\textbf{Artefacto 5: Práctica de Fundamentos (Set de Respaldo)}]
\end{tcolorbox}

\section*{Ejercicios Obligatorios (Independientes de la Ruta de Clase)[cite: 4]}

\textbf{Ejercicio 1 — Diagnóstico:}[cite: 4]
1. ¿Qué almacena un capacitor?[cite: 4]
2. ¿Qué almacena un inductor?[cite: 4]
3. ¿Qué significa el instante $0^+$?[cite: 4]
4. ¿Qué significa físicamente $\frac{dx}{dt}$?[cite: 4]

\textbf{Ejercicio 2 — Signos de las Derivadas:}[cite: 4]
Explique el significado físico de:
1. $\frac{dv_C}{dt} > 0$[cite: 4] \quad
2. $\frac{dv_C}{dt} = 0$[cite: 4] \quad
3. $\frac{di_L}{dt} < 0$[cite: 4]

\textbf{Ejercicio 3 — Relaciones Constitutivas:}[cite: 4]
\begin{itemize}
    \item \textbf{A:} Si $C=100\,\mu\text{F}$ e $i_C=1\,\text{mA}$, encuentre $\frac{dv_C}{dt}$. (Resultado esperado: $10\,\text{V/s}$)[cite: 4].
    \item \textbf{B:} Si $L=10\,\text{mH}$ y $v_L=-1\,\text{V}$, encuentre $\frac{di_L}{dt}$. (Resultado esperado: $-100\,\text{A/s}$)[cite: 4].
\end{itemize}

\textbf{Ejercicio 4 — Continuidad:}[cite: 4]
\begin{itemize}
    \item Si $v_C(0^-) = 5\,\text{V} \rightarrow$ halle $v_C(0^+)$.[cite: 4]
    \item Si $i_L(0^-) = 0.3\,\text{A} \rightarrow$ halle $i_L(0^+)$.[cite: 4]
\end{itemize}

\textbf{Ejercicio 5 — Ecuación RC Guiada:}[cite: 4]
Complete y luego explique en palabras qué modela la ecuación:
$$ V_s - \underline{\hspace{1.5cm}} - \underline{\hspace{1.5cm}} = 0 \quad \implies \quad \underline{\hspace{1.5cm}}\frac{dv_C}{dt} + v_C = \underline{\hspace{1.5cm}} $$[cite: 4]

\section*{Ejercicios Condicionales (Solo si se activó la Ruta A)[cite: 4]}

\textbf{Ejercicio 6:}[cite: 4]
Dado $R=3\,\text{k}\Omega$ y $C=100\,\mu\text{F}$, halle $\tau$. (Esperado: $0.3\,\text{s}$)[cite: 4]. Compare con un circuito de $R=1\,\text{k}\Omega$ y mismo $C$. ¿Cuál responde más rápido? (Esperado: el circuito de $1\,\text{k}\Omega$)[cite: 4].

\textbf{Ejercicio 7:}[cite: 4]
Dado $x(0^+) = 6$ y $x(\infty) = 2$[cite: 4]. Sin calcular la exponencial: 1) Identifique el valor de partida, 2) Identifique el valor final, 3) Indique si la respuesta sube o baja. (Esperado: Empieza en 6, termina en 2, baja)[cite: 4].

\end{document}
